%%%%%%%%%%%%%%%%%%%%%%%%%%%%%%%%%%%%%%%%%
% The Legrand Orange Book
% LaTeX Template
% Version 1.3 (21/8/13)
%
% This template has been downloaded from:
% http://www.LaTeXTemplates.com
%
% Original author:
% Mathias Legrand (legrand.mathias@gmail.com)
%
% License:
% CC BY-NC-SA 3.0 (http://creativecommons.org/licenses/by-nc-sa/3.0/)
%
% Compiling this template:
% This template uses biber for its bibliography and makeindex for its index.
% When you first open the template, compile it from the command line with the 
% commands below to make sure your LaTeX distribution is configured correctly:
%
% 1) pdflatex main
% 2) makeindex main.idx -s StyleInd.ist
% 3) biber main
% 4) pdflatex main x 2
%
% After this, when you wish to update the bibliography/index use the appropriate
% command above and make sure to compile with pdflatex several times 
% afterwards to propagate your changes to the document.
%
% This template also uses a number of packages which may need to be
% updated to the newest versions for the template to compile. It is strongly
% recommended you update your LaTeX distribution if you have any
% compilation errors.
%
% Important note:
% Chapter heading images should have a 2:1 width:height ratio,
% e.g. 920px width and 460px height.
%
%%%%%%%%%%%%%%%%%%%%%%%%%%%%%%%%%%%%%%%%%

%----------------------------------------------------------------------------------------
%	PACKAGES AND OTHER DOCUMENT CONFIGURATIONS
%----------------------------------------------------------------------------------------

\documentclass[11pt,fleqn]{book} % Default font size and left-justified equations

\usepackage[top=3cm,bottom=3cm,left=3.2cm,right=3.2cm,headsep=10pt,a4paper]{geometry} % Page margins

\usepackage{xcolor} % Required for specifying colors by name
\definecolor{ocre}{RGB}{243,102,25} % Define the orange color used for highlighting throughout the book

\newcommand{\degree}{\ensuremath{^\circ}}
\usepackage{wrapfig}

% Font Settings
\usepackage{avant} % Use the Avantgarde font for headings
%\usepackage{times} % Use the Times font for headings
\usepackage{mathptmx} % Use the Adobe Times Roman as the default text font together with math symbols from the Sym­bol, Chancery and Com­puter Modern fonts

\usepackage{microtype} % Slightly tweak font spacing for aesthetics
\usepackage[utf8]{inputenc} % Required for including letters with accents
\usepackage[T1]{fontenc} % Use 8-bit encoding that has 256 glyphs

% Bibliography
\usepackage[style=alphabetic,sorting=nyt,sortcites=true,autopunct=true,babel=hyphen,hyperref=true,abbreviate=false,backref=true,backend=biber]{biblatex}
\addbibresource{bibliography.bib} % BibTeX bibliography file
\defbibheading{bibempty}{}

% Index
\usepackage{calc} % For simpler calculation - used for spacing the index letter headings correctly
\usepackage{makeidx} % Required to make an index
\makeindex % Tells LaTeX to create the files required for indexing

%----------------------------------------------------------------------------------------

%----------------------------------------------------------------------------------------
%	VARIOUS REQUIRED PACKAGES
%----------------------------------------------------------------------------------------

\usepackage{titlesec} % Allows customization of titles

\usepackage{graphicx} % Required for including pictures
\graphicspath{{img/}} % Specifies the directory where pictures are stored

\usepackage{lipsum} % Inserts dummy text

\usepackage{tikz} % Required for drawing custom shapes

\usepackage[english]{babel} % English language/hyphenation

\usepackage{enumitem} % Customize lists
\setlist{nolistsep} % Reduce spacing between bullet points and numbered lists

\usepackage{booktabs} % Required for nicer horizontal rules in tables

\usepackage{eso-pic} % Required for specifying an image background in the title page

%----------------------------------------------------------------------------------------
%	MAIN TABLE OF CONTENTS
%----------------------------------------------------------------------------------------

\usepackage{titletoc} % Required for manipulating the table of contents

\contentsmargin{0cm} % Removes the default margin
% Chapter text styling
\titlecontents{chapter}[1.25cm] % Indentation
{\addvspace{15pt}\large\sffamily\bfseries} % Spacing and font options for chapters
{\color{ocre!60}\contentslabel[\Large\thecontentslabel]{1.25cm}\color{ocre}} % Chapter number
{}  
{\color{ocre!60}\normalsize\sffamily\bfseries\;\titlerule*[.5pc]{.}\;\thecontentspage} % Page number
% Section text styling
\titlecontents{section}[1.25cm] % Indentation
{\addvspace{5pt}\sffamily\bfseries} % Spacing and font options for sections
{\contentslabel[\thecontentslabel]{1.25cm}} % Section number
{}
{\sffamily\hfill\color{black}\thecontentspage} % Page number
[]
% Subsection text styling
\titlecontents{subsection}[1.25cm] % Indentation
{\addvspace{1pt}\sffamily\small} % Spacing and font options for subsections
{\contentslabel[\thecontentslabel]{1.25cm}} % Subsection number
{}
{\sffamily\;\titlerule*[.5pc]{.}\;\thecontentspage} % Page number
[] 

%----------------------------------------------------------------------------------------
%	MINI TABLE OF CONTENTS IN CHAPTER HEADS
%----------------------------------------------------------------------------------------

% Section text styling
\titlecontents{lsection}[0em] % Indendating
{\footnotesize\sffamily} % Font settings
{}
{}
{}

% Subsection text styling
\titlecontents{lsubsection}[.5em] % Indentation
{\normalfont\footnotesize\sffamily} % Font settings
{}
{}
{}
 
%----------------------------------------------------------------------------------------
%	PAGE HEADERS
%----------------------------------------------------------------------------------------

\usepackage{fancyhdr} % Required for header and footer configuration

\pagestyle{fancy}
\renewcommand{\chaptermark}[1]{\markboth{\sffamily\normalsize\bfseries #1}{}} % Chapter text font settings
\renewcommand{\sectionmark}[1]{\markright{\sffamily\normalsize\thesection\hspace{5pt}#1}{}} % Section text font settings
\fancyhf{} \fancyhead[LE,RO]{\sffamily\normalsize\thepage} % Font setting for the page number in the header
\fancyhead[LO]{\rightmark} % Print the nearest section name on the left side of odd pages
\fancyhead[RE]{\leftmark} % Print the current chapter name on the right side of even pages
\renewcommand{\headrulewidth}{0.5pt} % Width of the rule under the header
\addtolength{\headheight}{2.5pt} % Increase the spacing around the header slightly
\renewcommand{\footrulewidth}{0pt} % Removes the rule in the footer
\fancypagestyle{plain}{\fancyhead{}\renewcommand{\headrulewidth}{0pt}} % Style for when a plain pagestyle is specified

% Removes the header from odd empty pages at the end of chapters
\makeatletter
\renewcommand{\cleardoublepage}{
\clearpage\ifodd\c@page\else
\hbox{}
\vspace*{\fill}
\thispagestyle{empty}
\newpage
\fi}

%----------------------------------------------------------------------------------------
%	THEOREM STYLES
%----------------------------------------------------------------------------------------

\usepackage{amsmath,amsfonts,amssymb,amsthm} % For including math equations, theorems, symbols, etc

\newcommand{\intoo}[2]{\mathopen{]}#1\,;#2\mathclose{[}}
\newcommand{\ud}{\mathop{\mathrm{{}d}}\mathopen{}}
\newcommand{\intff}[2]{\mathopen{[}#1\,;#2\mathclose{]}}
\newtheorem{notation}{Notation}[chapter]

%%%%%%%%%%%%%%%%%%%%%%%%%%%%%%%%%%%%%%%%%%%%%%%%%%%%%%%%%%%%%%%%%%%%%%%%%%%
%%%%%%%%%%%%%%%%%%%% dedicated to boxed/framed environements %%%%%%%%%%%%%%
%%%%%%%%%%%%%%%%%%%%%%%%%%%%%%%%%%%%%%%%%%%%%%%%%%%%%%%%%%%%%%%%%%%%%%%%%%%
\newtheoremstyle{ocrenumbox}% % Theorem style name
{0pt}% Space above
{0pt}% Space below
{\normalfont}% % Body font
{}% Indent amount
{\small\bf\sffamily\color{ocre}}% % Theorem head font
{\;}% Punctuation after theorem head
{0.25em}% Space after theorem head
{\small\sffamily\color{ocre}\thmname{#1}\nobreakspace\thmnumber{\@ifnotempty{#1}{}\@upn{#2}}% Theorem text (e.g. Theorem 2.1)
\thmnote{\nobreakspace\the\thm@notefont\sffamily\bfseries\color{black}---\nobreakspace#3.}} % Optional theorem note
\renewcommand{\qedsymbol}{$\blacksquare$}% Optional qed square

\newtheoremstyle{blacknumex}% Theorem style name
{5pt}% Space above
{5pt}% Space below
{\normalfont}% Body font
{} % Indent amount
{\small\bf\sffamily}% Theorem head font
{\;}% Punctuation after theorem head
{0.25em}% Space after theorem head
{\small\sffamily{\tiny\ensuremath{\blacksquare}}\nobreakspace\thmname{#1}\nobreakspace\thmnumber{\@ifnotempty{#1}{}\@upn{#2}}% Theorem text (e.g. Theorem 2.1)
\thmnote{\nobreakspace\the\thm@notefont\sffamily\bfseries---\nobreakspace#3.}}% Optional theorem note

\newtheoremstyle{blacknumbox} % Theorem style name
{0pt}% Space above
{0pt}% Space below
{\normalfont}% Body font
{}% Indent amount
{\small\bf\sffamily}% Theorem head font
{\;}% Punctuation after theorem head
{0.25em}% Space after theorem head
{\small\sffamily\thmname{#1}\nobreakspace\thmnumber{\@ifnotempty{#1}{}\@upn{#2}}% Theorem text (e.g. Theorem 2.1)
\thmnote{\nobreakspace\the\thm@notefont\sffamily\bfseries---\nobreakspace#3.}}% Optional theorem note

%%%%%%%%%%%%%%%%%%%%%%%%%%%%%%%%%%%%%%%%%%%%%%%%%%%%%%%%%%%%%%%%%%%%%%%%%%%
%%%%%%%%%%%%% dedicated to non-boxed/non-framed environements %%%%%%%%%%%%%
%%%%%%%%%%%%%%%%%%%%%%%%%%%%%%%%%%%%%%%%%%%%%%%%%%%%%%%%%%%%%%%%%%%%%%%%%%%
\newtheoremstyle{ocrenum}% % Theorem style name
{5pt}% Space above
{5pt}% Space below
{\normalfont}% % Body font
{}% Indent amount
{\small\bf\sffamily\color{ocre}}% % Theorem head font
{\;}% Punctuation after theorem head
{0.25em}% Space after theorem head
{\small\sffamily\color{ocre}\thmname{#1}\nobreakspace\thmnumber{\@ifnotempty{#1}{}\@upn{#2}}% Theorem text (e.g. Theorem 2.1)
\thmnote{\nobreakspace\the\thm@notefont\sffamily\bfseries\color{black}---\nobreakspace#3.}} % Optional theorem note
\renewcommand{\qedsymbol}{$\blacksquare$}% Optional qed square
\makeatother

% Defines the theorem text style for each type of theorem to one of the three styles above
\newcounter{dummy} 
\numberwithin{dummy}{section}
\theoremstyle{ocrenumbox}
\newtheorem{theoremeT}[dummy]{Theorem}
\newtheorem{problem}{Problem}[chapter]
\newtheorem{exerciseT}{Exercise}[chapter]
\theoremstyle{blacknumex}
\newtheorem{exampleT}{Example}[chapter]
\theoremstyle{blacknumbox}
\newtheorem{vocabulary}{Vocabulary}[chapter]
\newtheorem{definitionT}{Definition}[section]
\newtheorem{corollaryT}[dummy]{Corollary}
\theoremstyle{ocrenum}
\newtheorem{proposition}[dummy]{Proposition}

%----------------------------------------------------------------------------------------
%	DEFINITION OF COLORED BOXES
%----------------------------------------------------------------------------------------

\RequirePackage[framemethod=default]{mdframed} % Required for creating the theorem, definition, exercise and corollary boxes

% Theorem box
\newmdenv[skipabove=7pt,
skipbelow=7pt,
backgroundcolor=black!5,
linecolor=ocre,
innerleftmargin=5pt,
innerrightmargin=5pt,
innertopmargin=5pt,
leftmargin=0cm,
rightmargin=0cm,
innerbottommargin=5pt]{tBox}

% Exercise box	  
\newmdenv[skipabove=7pt,
skipbelow=7pt,
rightline=false,
leftline=true,
topline=false,
bottomline=false,
backgroundcolor=ocre!10,
linecolor=ocre,
innerleftmargin=5pt,
innerrightmargin=5pt,
innertopmargin=5pt,
innerbottommargin=5pt,
leftmargin=0cm,
rightmargin=0cm,
linewidth=4pt]{eBox}	

% Definition box
\newmdenv[skipabove=7pt,
skipbelow=7pt,
rightline=false,
leftline=true,
topline=false,
bottomline=false,
linecolor=ocre,
innerleftmargin=5pt,
innerrightmargin=5pt,
innertopmargin=0pt,
leftmargin=0cm,
rightmargin=0cm,
linewidth=4pt,
innerbottommargin=0pt]{dBox}	

% Corollary box
\newmdenv[skipabove=7pt,
skipbelow=7pt,
rightline=false,
leftline=true,
topline=false,
bottomline=false,
linecolor=gray,
backgroundcolor=black!5,
innerleftmargin=5pt,
innerrightmargin=5pt,
innertopmargin=5pt,
leftmargin=0cm,
rightmargin=0cm,
linewidth=4pt,
innerbottommargin=5pt]{cBox}				  
		  

% Creates an environment for each type of theorem and assigns it a theorem text style from the "Theorem Styles" section above and a colored box from above
\newenvironment{theorem}{\begin{tBox}\begin{theoremeT}}{\end{theoremeT}\end{tBox}}
\newenvironment{exercise}{\begin{eBox}\begin{exerciseT}}{\hfill{\color{ocre}\tiny\ensuremath{\blacksquare}}\end{exerciseT}\end{eBox}}				  
\newenvironment{definition}{\begin{dBox}\begin{definitionT}}{\end{definitionT}\end{dBox}}	
\newenvironment{example}{\begin{exampleT}}{\hfill{\tiny\ensuremath{\blacksquare}}\end{exampleT}}		
\newenvironment{corollary}{\begin{cBox}\begin{corollaryT}}{\end{corollaryT}\end{cBox}}	

%----------------------------------------------------------------------------------------
%	REMARK ENVIRONMENT
%----------------------------------------------------------------------------------------

\newenvironment{remark}{\par\vskip10pt\small % Vertical white space above the remark and smaller font size
\begin{list}{}{
\leftmargin=35pt % Indentation on the left
\rightmargin=25pt}\item\ignorespaces % Indentation on the right
\makebox[-2.5pt]{\begin{tikzpicture}[overlay]
\node[draw=ocre!60,line width=1pt,circle,fill=ocre!25,font=\sffamily\bfseries,inner sep=2pt,outer sep=0pt] at (-15pt,0pt){\textcolor{ocre}{R}};\end{tikzpicture}} % Orange R in a circle
\advance\baselineskip -1pt}{\end{list}\vskip5pt} % Tighter line spacing and white space after remark

%----------------------------------------------------------------------------------------
%	SECTION NUMBERING IN THE MARGIN
%----------------------------------------------------------------------------------------

\makeatletter
\renewcommand{\@seccntformat}[1]{\llap{\textcolor{ocre}{\csname the#1\endcsname}\hspace{1em}}}                    
\renewcommand{\section}{\@startsection{section}{1}{\z@}
{-4ex \@plus -1ex \@minus -.4ex}
{1ex \@plus.2ex }
{\normalfont\large\sffamily\bfseries}}
\renewcommand{\subsection}{\@startsection {subsection}{2}{\z@}
{-3ex \@plus -0.1ex \@minus -.4ex}
{0.5ex \@plus.2ex }
{\normalfont\sffamily\bfseries}}
\renewcommand{\subsubsection}{\@startsection {subsubsection}{3}{\z@}
{-2ex \@plus -0.1ex \@minus -.2ex}
{0.2ex \@plus.2ex }
{\normalfont\small\sffamily\bfseries}}                        
\renewcommand\paragraph{\@startsection{paragraph}{4}{\z@}
{-2ex \@plus-.2ex \@minus .2ex}
{0.1ex}
{\normalfont\small\sffamily\bfseries}}

%----------------------------------------------------------------------------------------
%	CHAPTER HEADINGS
%----------------------------------------------------------------------------------------

\newcommand{\thechapterimage}{}
\newcommand{\chapterimage}[1]{\renewcommand{\thechapterimage}{#1}}
\def\thechapter{\arabic{chapter}}
\def\@makechapterhead#1{
\thispagestyle{empty}
{\centering \normalfont\sffamily
\ifnum \c@secnumdepth >\m@ne
\if@mainmatter
\startcontents
\begin{tikzpicture}[remember picture,overlay]
\node at (current page.north west)
{\begin{tikzpicture}[remember picture,overlay]

\node[anchor=north west,inner sep=0pt] at (0,0) {\includegraphics[width=\paperwidth]{\thechapterimage}};

%Commenting the 3 lines below removes the small contents box in the chapter heading
\draw[fill=white,opacity=.6] (1cm,0) rectangle (8cm,-7cm);
\node[anchor=north west] at (1cm,.25cm) {\parbox[t][8cm][t]{6.5cm}{\huge\bfseries\flushleft \printcontents{l}{1}{\setcounter{tocdepth}{2}}}};

\draw[anchor=west] (5cm,-9cm) node [rounded corners=25pt,fill=white,fill opacity=.6,text opacity=1,draw=ocre,draw opacity=1,line width=2pt,inner sep=15pt]{\huge\sffamily\bfseries\textcolor{black}{\thechapter\ ---\ #1\vphantom{plPQq}\makebox[22cm]{}}};
\end{tikzpicture}};
\end{tikzpicture}}\par\vspace*{230\p@}
\fi
\fi
}
\def\@makeschapterhead#1{
\thispagestyle{empty}
{\centering \normalfont\sffamily
\ifnum \c@secnumdepth >\m@ne
\if@mainmatter
\startcontents
\begin{tikzpicture}[remember picture,overlay]
\node at (current page.north west)
{\begin{tikzpicture}[remember picture,overlay]
\node[anchor=north west] at (-4pt,4pt) {\includegraphics[width=\paperwidth]{\thechapterimage}};
\draw[anchor=west] (5cm,-9cm) node [rounded corners=25pt,fill=white,opacity=.7,inner sep=15.5pt]{\huge\sffamily\bfseries\textcolor{black}{\vphantom{plPQq}\makebox[22cm]{}}};
\draw[anchor=west] (5cm,-9cm) node [rounded corners=25pt,draw=ocre,line width=2pt,inner sep=15pt]{\huge\sffamily\bfseries\textcolor{black}{#1\vphantom{plPQq}\makebox[22cm]{}}};
\end{tikzpicture}};
\end{tikzpicture}}\par\vspace*{230\p@}
\fi
\fi
}
\makeatother
 % Insert the commands.tex file which contains the majority of the structure behind the template

\begin{document}

%----------------------------------------------------------------------------------------
%	TITLE PAGE
%----------------------------------------------------------------------------------------

\begingroup
\thispagestyle{empty}
\AddToShipoutPicture*{\put(6,5){\includegraphics[scale=1]{background}}} % Image background
\centering
\vspace*{9cm}
\par\normalfont\fontsize{35}{35}\sffamily\selectfont
\Huge Brew Master 9000\\
Hardware\\
\textit{Research and reference}\par % Book title
\vspace*{1cm}
{\Large Iver Egge\\
Kristoffer Dalby\\
Johan Slettvold\\
Christer Lund\\} %\par % Author name
\endgroup

%----------------------------------------------------------------------------------------
%	COPYRIGHT PAGE
%----------------------------------------------------------------------------------------

\newpage
~\vfill
\thispagestyle{empty}

\noindent Copyright \copyright\ 2013 Webrewthebestbeer1\\ % Copyright notice

\noindent \textsc{Published by null}\\ % Publisher

\noindent \textsc{wb3.no}\\ % URL

\noindent We don't need no license\\ % License information

\noindent The U.S. customary units should have been aborted\\

\noindent \textit{First printing, October 2013} % Printing/edition date

%----------------------------------------------------------------------------------------
%	TABLE OF CONTENTS
%----------------------------------------------------------------------------------------

\chapterimage{chapter_head_1.pdf} % Table of contents heading image

\pagestyle{empty} % No headers

\tableofcontents % Print the table of contents itself

\cleardoublepage % Forces the first chapter to start on an odd page so it's on the right

\pagestyle{fancy} % Print headers again

%----------------------------------------------------------------------------------------
%	PART 1
%----------------------------------------------------------------------------------------

\part{Theory}

%----------------------------------------------------------------------------------------

\chapterimage{tubes_fittings.pdf} % Chapter heading image

\chapter{Tubes and fittings}

\section{Tri Clover / Tri Clamp / Sanitary}\index{Tri Clover / Tri Clamp / Sanitary}

These fittings are used throughout the brewing system as they provide easy dissassembling and just look really, really cool.\\

They consist of two flanges and a gasket that are compressed together with a clamp. The clamp is tightened with a thumb screw.

\subsection{Gaskets}\index{Tri Clover / Tri Clamp / Sanitary!Gaskets}

There are many types of gaskets produced for use with tri clamp fittings. Three of these types are commonly used in amateur brewing systems; silocone, EPDM (ethylene propylene diene monomer) rubber and PTFE (polytetrafluoroethylene) teflon.

\subsubsection{Silicone}

Temperature rating of $-49\degree$C \textasciitilde $230\degree$C.\\
These are the most common gaskets that are available. They will degrade over time if used with strong acids. As they stick very well to metal and are soft, they provide an excellent seal.

\subsubsection{EPDM rubber}

Temperature rating of $-34\degree$C \textasciitilde $149\degree$C.\\
They have better chemical resistance than silicone gaskets and will last longer than silicone gaskets. As the name implies they are made of rubber and are therefor soft and somewhat sticky.

\subsubsection{PTFE teflon}

Temperature rating of $-73\degree$C \textasciitilde $260\degree$C.\\
They have the best chemical resistance of all gaskets and will last the longest. They are, however, hard and will need considerably more compression to provide a good seal.

\subsubsection{Gaskets with flanges}

A normal gasket will fall right of the fitting when loose. If the gasket has a flange that covers the outer part of the fitting it will stay on the fitting when dissassembling (i.e. does not fall into the warm wort).

\subsubsection{Stiffness}

A stiff gasket that is not sticky will allow you to turn the fitting without dissassembling the entire connection. Stiffer gaskets will need more compression to provide a good seal.

\section{Compression fittings}\index{Compression fittings}

Compression fittings consists of a compression nut and ring that slides over the tube and a threaded fitting. If the tube is made of soft metal there should also be a support insert that is inserted into the tube to prevent it from collapsing.

These fittings should not be over tightened as this will ruin the compression ring and therefore the seal.

\section{Valves}\index{Valves}

Any normal ball valve will work in a brewing system. Only full port ball valves have the same size hole in the ball as the pipeline.

It should be possible to dissassemble the valves for cleaning and repair. Therefore 3-part valves are recommended for a brewing system.

%----------------------------------------------------------------------------------------
%	CHAPTER 2
%----------------------------------------------------------------------------------------

\chapterimage{pumps.pdf} % Chapter heading image

\chapter{Pump}

\section{Requirements}\index{Pump!Requirements}

A suitable pump for a brewing system have the following requirements

\begin{itemize}
\item All parts have to be of food grade.
\item It should be magnetic coupled so that in the event of the impeller becoming stuck due to malt particles, the motor will not burn out.
\item The lift limit should be at least 2 meters.
\item The temperature rating should be at least that of boiling wort ($100\degree$C).
\item It should be self-priming.
\end{itemize}

\section{Priming}\index{Pump!Priming}

Priming is to fill the pump head with the liquid that is to be pumped. As the pump is stored dry, the contents of the pump head is only air. As air and water have very different physical properties, a pump that is designed to pump a liquid will perform terrible at pumping air and will eventually break.\\

A self-priming pump differs from a non-priming pump that it can pump a mixture of air and liquid. This means that the pump will be able to remove air trapped in the head as long as it has a source of liquid.\\

A pump made to pump liquid should NEVER be run dry.

\section{Mounting}\index{Pump!Mounting}

The pump should be mounted so that liquid enters at the bottom of the head and exits at the top. This is to prevent air pockets from occupying the head. If the pump has to be mounted in a vertical position the head should be placed at the top, not the bottom.\\

Many pumps require that there is pressure at the inlet. Therefore it should be mounted as far as possible below the source of liquid.

\section{Cavitation}

Cavitation is the formation of vapor cavities in a liquid due to pressure differences around the impeller as it turns quickly and with great force. It can be harmful to the pump and will produce a rattle-like sound.

%----------------------------------------------------------------------------------------

\chapterimage{temp.pdf} % Chapter heading image

\chapter{Sensors and input}

\section{Temperature measurements}

\subsection{Probes}\index{Probes}

There are many different temperature probes available that either communicates digitally or outputs an analog signal that can be processed for use with various microcontrollers. The most common types seem to be

\begin{itemize}
\item DS18B20 (digital, Dallas 1-wire)
\item LM35 (analog, 0-1V)
\item Thermocouple (analog, needs amplifier)
\end{itemize}

The DS18B20 and LM35 are quite easy to use but lacks good casings for wet use. Thermocouples are the industry standard for temperature measurements and comes in a great number of different casings. They consist of two different types of conductors that provide a voltage difference according to the temperature of the system. This voltage difference is in the order of microvolts and an amplifier is needed to read the output.

\subsubsection{DS18B20}

The DS18B20 comes in a standarized TO-92 packaging. It communicates over the Dallas 1-wire bus that enables several sensors to share a dataline. They can also source power from the dataline so that only two wires are needed for the entire sensor-array.
The temperature range of this sensor is $-55\degree$C \textasciitilde $125\degree$C and it provides an accuracy of $\pm$ $0.5\degree$C in the $-10\degree$C \textasciitilde $85\degree$C range.

\subsubsection{Thermocouples}

Thermocouples consists of two conductors that differ in physical properties. When these two conductors contact each other they produce a voltage proportional to the temperature difference.

The main dissadvantage of thermocouples is that a system error of less than one degree Kelvin is difficult to achieve.

\subsubsection{K-type thermocouple and amplifier}

This seems to be the most common probe of the thermocouples types. It consists of chromel-alumel conductors. The temperature range of the probe can be as wide as $-200\degree$C \textasciitilde $1350\degree$C.

The MAX31855 is an amplifier for thermocouples produced by Maxim Integrated. It outputs an analog signal and requires an ADC for use with a microcontroller (or a microcontroller with an integrated ADC). There are libraries available for the Ardunio and the Raspberry Pi. 

\subsection{Placement}

Liquid that is being heated does not have uniformly distributed temperature. There should be two probes per tank, and the average of the two will be the actual temperature of the system.
If only one probe is used to measure the temperature of a tank the probe should be fitted at the middle (why?).

\section{Flow sensors / Flow meter}

A flow sensor can essensially be as simple as a fan that is rotated by the moving liquid. A magnet on the fan axel moves past a magnetic sensor and produces an electronic pulse on the output line. A microcontroller then records these pulses and is able to calculate the volume of liquid that has passed through the sensor.

These simple flow sensors comes in different volume per time ratings. 

%----------------------------------------------------------------------------------------

\chapterimage{relays.pdf} % Chapter heading image

\chapter{Relays and heatsinks}

\section{Electromechanical relays}

Electromechanical relays operate by an electromagnet that mechanically moves a conductor.

These need a flyback diode in parallel with the coil to prevent the controlling circuitry from damage when the relay is turned off and the magnetic field stored in the coil collapses. 

\section{Solid state relays}

\subsection{Zero crossing}

Zero crossing means that the relay switches from the conducting to the non-conducting phase when the AC sine wave reaches the zero crossing point. This reduces the surge current through the load and greatly improves the lifespan of the SSR.

\subsection{Calculating heatsink requirements for SSR}

SSR's generally produce about 1W of heat for every ampere of the load.

\subsection{Thermal pads and compound}

Thermal pads or thermal compound should always be used when connecting a SSR to a heatsink.

As tiny irregularities in the flat, metal surface of the SSR and heatsink provide trapped air and thus insulation, thermal pads or compound is used to fill in these irregularities. If compound is used, only a very thin film is necessary to provide good thermal conductivity.

%----------------------------------------------------------------------------------------

\chapterimage{rims.pdf} % Chapter heading image

\chapter{RIMS-tube}

\section{What is RIMS}

RIMS stands for Recirculating Infusion Mash System.\\
The wort is recirculated continously through the malt. The RIMS-tube contains a heating element and a temperature probe that enables the brewer to do step mashing by pumping the wort through the tube. RIMS also provides crystal clear wort as it is filtered continously through the malt.

\section{Components}

A typical RIMS-tube consists of the following

\begin{itemize}
\item Two tee's
\item A heating element adapter and heating element
\item A temperature probe adapater and temperature probe
\item Two hose/pipe connections
\end{itemize}

%----------------------------------------------------------------------------------------

\chapterimage{vessel.pdf} % Chapter heading image

\chapter{Tanks, vessels and kettles}

\section{Volumes}

The following equations for volumes do not take into consideration other objects in the tanks, such as filters and stirrers. They instead give the minimum volume of the containers. Headspace and other objects should also be included in the final volume of the tanks.

\subsection{Hot Liquor Tun (HLT)}

This tank provides hot water for the mash and sparge. It has a heating element installed in the tank and uses the RIMS-tube for additional power when heating the mash water.

It's volume should be at least that of the needed mash water for the brew. This is calculated as follows

\begin{equation}
V_{HLT} = m_{malt} \cdot d_{malt} \cdot {379 \over 400}
%\label{Volume of mash water based on weight of malt and malt thickness}
\end{equation}

Where d$_{malt}$ is the thickness of the malt and is generally set to 2.61 $L\over kg $.

\subsection{Mash/Lauter Tun (MLT)}

This tank is where the malt is mashed at a constant or several, precise temperatures (step mashing). To keep the temperature constant, or to rise the temperature, the wort is passed through the RIMS-tube which heats the wort. The tank should also be very well insulated.

It's volume is calculated as follows

\begin{equation}
V_{MLT} = m_{malt} \cdot (0.67 + d_{malt})
\end{equation}

\subsection{Boil Kettle (BK)}

This tank is where the wort is boiled with the addition of hops. It has two heating elements installed in the tank and has an outlet to a chiller.

It's volume should be that of the wort produced for pre-boil and as a rule of thumb have a 20\% head room to compensate for foam and splash.

\begin{equation}
V_{wort, pre-boil} = {{V_{beer} + V_{hops, absorbed}} \over 0.96} \cdot {1 \over k_{evap}}
\end{equation}

Where k$_{evap}$ is the evaporisationcoefficient which is based on the power from the heating elements and the surface area of the top of the kettle. It is generally set to 0.95. V$_{hops, absorbed}$ is the volume of wort that is absorbed by the hops and is dependent on the amount of hops. 100g of hops will probably absorb 1L of wort.

The total volume of the boil kettle should then be

\begin{equation}
V_{BK} = V_{wort, pre-boil} \cdot 1.2
\end{equation}

\section{Insulation}

The formula for calculating heat loss based on insulation is

\begin{equation}
H = {{A (T_h - T_c)} \over R}
\end{equation}

Where A is the surface area of the tank, T$_h$ is the temperature of the liquid in the tank and T$_c$ is the temperature of the surroundings. R is a unique value for the insulation material, where a higher R-value designates better insulation properties.\\

Many insulation materials are based on trapping air in the material. Air is a very bad conductor of heat, and therefore the material insulates heat transfer. These insulation materials perform worse if they are compressed, as they can not trap enough air.


%----------------------------------------------------------------------------------------

\chapterimage{mill.pdf} % Chapter heading image

\chapter{Malt mills}

\section{Volumes}

%----------------------------------------------------------------------------------------

\chapterimage{equations.pdf} % Chapter heading image

\chapter{Misc. equations}

\section{Heat}

\subsection{Newtons law of cool'}

\subsection{Heat capacity}

\section{Electricity}

\subsection{Ohms law}

%----------------------------------------------------------------------------------------
%	PART 2
%----------------------------------------------------------------------------------------

\part{Practice}

%----------------------------------------------------------------------------------------
%	PART 2
%----------------------------------------------------------------------------------------

\part{Bibliography and index}

%----------------------------------------------------------------------------------------
%	BIBLIOGRAPHY
%----------------------------------------------------------------------------------------

\chapterimage{chapter_head_1.pdf} % Chapter heading image

\chapter*{Bibliography}
\addcontentsline{toc}{chapter}{\textcolor{ocre}{Bibliography}}
\section*{Books}
\addcontentsline{toc}{section}{Books}
\printbibliography[heading=bibempty,type=book]
\section*{Articles}
\addcontentsline{toc}{section}{Articles}
\printbibliography[heading=bibempty,type=article]

%----------------------------------------------------------------------------------------
%	INDEX
%----------------------------------------------------------------------------------------

\cleardoublepage
\setlength{\columnsep}{0.75cm}
\addcontentsline{toc}{chapter}{\textcolor{ocre}{Index}}
\printindex

%----------------------------------------------------------------------------------------

\end{document}